\documentclass[12pt, a4paper]{article}

\usepackage[utf8]{inputenc}
\usepackage[T1]{fontenc}
\usepackage[italian]{babel}
\usepackage{geometry}
\usepackage{graphicx}
\usepackage{amsmath}
\usepackage{amssymb}
\usepackage{listings}
\usepackage{xcolor}
\usepackage{hyperref}
\usepackage{booktabs}
\usepackage{float}
\usepackage{fancyhdr}
\usepackage{enumitem}
\usepackage[expansion=false]{microtype}
\usepackage{parskip}
\usepackage{setspace}

\geometry{
    top=2.5cm,
    bottom=2.5cm,
    left=3cm,
    right=3cm
}

\pagestyle{fancy}
\fancyhf{}
\rhead{TLN --- Lab 4}
\lhead{Università degli Studi di Torino}
\rfoot{\thepage}
\renewcommand{\headrulewidth}{0.4pt}

\definecolor{codegreen}{rgb}{0,0.6,0}
\definecolor{codegray}{rgb}{0.5,0.5,0.5}
\definecolor{codepurple}{rgb}{0.58,0,0.82}
\definecolor{backcolour}{rgb}{0.97,0.97,0.97}
\definecolor{codered}{rgb}{0.8,0,0}
\definecolor{darkblue}{rgb}{0.0,0.2,0.5}

\lstdefinestyle{pythonstyle}{
    backgroundcolor=\color{backcolour},
    commentstyle=\color{codegreen}\itshape,
    keywordstyle=\color{codepurple}\bfseries,
    numberstyle=\tiny\color{codegray},
    stringstyle=\color{codered},
    basicstyle=\ttfamily\small,
    breakatwhitespace=false,
    breaklines=true,
    captionpos=b,
    keepspaces=true,
    numbers=left,
    numbersep=8pt,
    showspaces=false,
    showstringspaces=false,
    showtabs=false,
    tabsize=2,
    frame=single,
    rulecolor=\color{codegray},
    xleftmargin=15pt
}
\lstset{style=pythonstyle}

\hypersetup{
    colorlinks=true,
    linkcolor=darkblue,
    filecolor=magenta,
    urlcolor=darkblue,
}

\onehalfspacing

% ============================================================
\begin{document}
% ============================================================

\begin{titlepage}
    \centering
    \vspace*{0.8cm}

    {\Large \textbf{Università degli Studi di Torino}}\\[0.2cm]
    {\large Dipartimento di Informatica}\\[0.1cm]
    {\large Corso di Laurea Magistrale in Informatica}

    \vspace{1.5cm}
    \rule{\linewidth}{0.8mm}\\[0.6cm]

    {\LARGE \textbf{Tecnologie del Linguaggio Naturale}}\\[0.4cm]
    {\Large \textbf{Lab 4 --- Modellazione del Linguaggio con N-gram}}\\[0.3cm]
    {\large \textit{Modeling social media and literary language with N-grams}}\\[0.6cm]

    \rule{\linewidth}{0.8mm}

    \vspace{1.8cm}

    {\large \textbf{Gruppo di lavoro:}}\\[0.5cm]

    \begin{tabular}{ll}
        \textbf{Studente}       & \textbf{Matricola} \\[0.3cm]
        Alessandro Olivero      & 915069  \\[0.2cm]
        Samuele Iudice          & 1235245 \\[0.2cm]
        Andrea Franco           & 1226739 \\[0.5cm]
        \textbf{Docente:}       & Prof. Daniele Radicioni \\[0.2cm]
        \textbf{Anno Accademico:} & 2024/2025 \\
    \end{tabular}

    \vfill
    {\large Torino, \today}
\end{titlepage}

\tableofcontents
\newpage

% ============================================================
\section{Introduzione}
% ============================================================

L'obiettivo di questo laboratorio è modellare e confrontare due varietà
linguistiche radicalmente diverse attraverso modelli N-gram: il linguaggio
dei \textbf{social media} (Twitter) e il \textbf{linguaggio letterario}
ottocentesco (\textit{Emma} di Jane Austen, 1816).

L'intuizione di fondo è semplice: un tweet e una frase di un romanzo dell'Ottocento
sono testi profondamente diversi, e questa diversità dovrebbe essere catturabile
in modo quantitativo. I modelli N-gram, pur essendo tra gli strumenti più semplici
del NLP, si prestano bene a questo tipo di analisi: contano le sequenze di parole
più frequenti e ne stimano le probabilità, rivelando le strutture ricorrenti
tipiche di ciascun dominio.

\subsection{Obiettivi specifici}

\begin{itemize}[leftmargin=1.5cm]
    \item Estrarre e confrontare gli N-gram (unigrammi, bigrammi, trigrammi)
          più frequenti nei due corpora, con e senza stopwords.
    \item Addestrare modelli di linguaggio MLE e Laplace e mostrarne
          differenze e limiti pratici.
    \item Analizzare la \textit{perplexity} cross-domain e discuterne i limiti
          quando i vocabolari hanno dimensioni diverse.
    \item Implementare un classificatore basato su \textit{log-likelihood}
          che superi i limiti della perplexity.
    \item Quantificare le differenze linguistiche attraverso il
          Type-Token Ratio (TTR) e la lunghezza media delle frasi.
    \item Generare testo nei due stili e confrontare qualitativamente
          i risultati.
\end{itemize}

\subsection{Strumenti e librerie}

\begin{itemize}[leftmargin=1.5cm]
    \item \textbf{Python 3.x} — linguaggio di implementazione
    \item \textbf{NLTK} --- tokenizzazione, corpora, modelli N-gram
    \item \textbf{Matplotlib} --- visualizzazione dei risultati
    \item \textbf{re} (regex) --- pulizia dei testi
\end{itemize}

% ============================================================
\section{Dataset}
% ============================================================

\subsection{Twitter Samples}

Il corpus Twitter proviene dal pacchetto \texttt{nltk.corpus.twitter\_samples},
che include tweet annotati automaticamente in base alla presenza di emoticon
positive o negative, raccolti nel 2015.

\begin{itemize}[leftmargin=1.5cm]
    \item \texttt{positive\_tweets.json} --- 5.000 tweet con emoticon \texttt{:)}
    \item \texttt{negative\_tweets.json} --- 5.000 tweet con emoticon \texttt{:(}
    \item \textbf{Totale}: 10.000 tweet
\end{itemize}

Si è scelto di usare questo sottoinsieme invece del file
\texttt{tweets.20150430-223406.json} (20.000 tweet generici),
perché quest'ultimo era dominato da contenuti politici relativi
alle elezioni britanniche del 30 aprile 2015 (menzionava Ed Miliband,
David Cameron, SNP), rendendo il confronto meno significativo.
I tweet positivi/negativi hanno invece un linguaggio più vario e
rappresentativo del social media in generale.

\begin{lstlisting}[language=Python, caption=Caricamento del corpus Twitter]
from nltk.corpus import twitter_samples

tweets_pos = twitter_samples.strings('positive_tweets.json')
tweets_neg = twitter_samples.strings('negative_tweets.json')
tweets_raw = tweets_pos + tweets_neg  # 10.000 tweet totali
\end{lstlisting}

Esempi rappresentativi:
\begin{itemize}[leftmargin=1.5cm]
    \item \textit{Positivo}: \texttt{\#FollowFriday @France\_Inte for being top engaged members this week :)}
    \item \textit{Negativo}: \texttt{hopeless for tmr :(}
\end{itemize}

\subsection{Corpus Letterario --- Emma di Jane Austen}

Il testo letterario è stato estratto dal corpus Gutenberg di NLTK,
che raccoglie opere classiche di dominio pubblico. È stato scelto
\textbf{\textit{Emma} di Jane Austen (1816)} come rappresentante
del linguaggio letterario formale ottocentesco: il romanzo è lungo,
narrativamente ricco, e utilizza costruzioni sintattiche complesse
con un vocabolario formale e preciso.

\begin{lstlisting}[language=Python, caption=Caricamento del corpus letterario]
from nltk.corpus import gutenberg

testo_letterario_raw = gutenberg.raw('austen-emma.txt')
# 887.071 caratteri totali
\end{lstlisting}

\subsection{Statistiche comparative}

La tabella seguente riassume le principali caratteristiche dei due corpora
dopo il preprocessing.

\begin{table}[H]
\centering
\begin{tabular}{lrr}
\toprule
\textbf{Metrica} & \textbf{Twitter} & \textbf{Letterario} \\
\midrule
Numero di testi      & 10.000 tweet & 1 romanzo \\
Token totali         & 94.706       & 158.267 \\
Parole uniche        & 12.024       & 9.305 \\
Lunghezza media      & 16,0 parole  & 25,6 parole \\
Type-Token Ratio     & 0,1270       & 0,0588 \\
\bottomrule
\end{tabular}
\caption{Statistiche comparative dei due corpora dopo preprocessing}
\end{table}

% ============================================================
\section{Preprocessing}
% ============================================================

Prima di costruire i modelli N-gram, i testi devono essere puliti
e normalizzati. I due corpora richiedono trattamenti diversi.

\subsection{Pulizia dei tweet}

I tweet contengono elementi non linguistici che andrebbero a
sporcare il modello: URL, menzioni utente (\texttt{@user}),
simboli degli hashtag, emoticon e punteggiatura.

\begin{lstlisting}[language=Python, caption=Funzione di pulizia dei tweet]
def pulisci_tweet(testo):
    testo = re.sub(r'http\S+', '', testo)      # rimuove URL
    testo = re.sub(r'@\w+', '', testo)          # rimuove @menzioni
    testo = re.sub(r'#', '', testo)             # rimuove simbolo #
    testo = re.sub(r'[^a-zA-Z\s]', '', testo)  # solo lettere
    testo = testo.lower().strip()
    return testo
\end{lstlisting}

Esempio di trasformazione:
\begin{itemize}[leftmargin=1.5cm]
    \item \textbf{Prima:} \texttt{\#FollowFriday @France\_Inte for being top engaged :)}
    \item \textbf{Dopo:}  \texttt{followfriday for being top engaged}
\end{itemize}

\subsection{Pulizia del testo letterario}

Per il romanzo è sufficiente rimuovere la punteggiatura e normalizzare
il caso, dato che non contiene URL o menzioni.

\begin{lstlisting}[language=Python, caption=Funzione di pulizia del testo letterario]
def pulisci_letterario(testo):
    testo = re.sub(r'[^a-zA-Z\s]', '', testo)
    testo = testo.lower().strip()
    return testo
\end{lstlisting}

\subsection{Tokenizzazione}

La tokenizzazione converte la stringa di testo in una lista di token
(parole). Si è usata la funzione \texttt{word\_tokenize} di NLTK
invece del semplice \texttt{split()}, perché gestisce correttamente
le contrazioni inglesi e separa la punteggiatura residua.

\begin{table}[H]
\centering
\begin{tabular}{ll}
\toprule
\textbf{Metodo} & \textbf{Output su ``I can't stop,''} \\
\midrule
\texttt{split()} & \texttt{["I", "can't", "stop,"]} \\
\texttt{word\_tokenize()} & \texttt{["I", "ca", "n't", "stop", ","]} \\
\bottomrule
\end{tabular}
\caption{Differenza tra split() e word\_tokenize()}
\end{table}

\subsection{Rimozione delle stopwords}

Le stopwords sono parole molto frequenti in qualsiasi testo inglese
(``the'', ``a'', ``in'', ``of''...) e non sono discriminative tra
i due domini. La loro rimozione permette di far emergere i bigrammi
davvero caratteristici di ciascun corpus.

\begin{lstlisting}[language=Python, caption=Rimozione delle stopwords]
from nltk.corpus import stopwords

stop = set(stopwords.words('english'))
tokens_tweet_filtrati = [t for t in tokens_tweet if t not in stop]
tokens_lett_filtrati  = [t for t in tokens_lett  if t not in stop]
\end{lstlisting}

% ============================================================
\section{Analisi N-gram}
% ============================================================

\subsection{Definizione}

Un \textbf{N-gram} è una sequenza contigua di $N$ parole in un testo.
Dato un testo $w_1, w_2, \ldots, w_T$, gli N-gram di ordine $N$ sono
tutte le sequenze $(w_i, w_{i+1}, \ldots, w_{i+N-1})$ per
$i = 1, \ldots, T-N+1$.

\begin{itemize}[leftmargin=1.5cm]
    \item \textbf{Unigrammi} ($N=1$): parole singole.
          Es.: ``i love cats'' $\to$ \texttt{[i, love, cats]}
    \item \textbf{Bigrammi} ($N=2$): coppie consecutive.
          Es.: ``i love cats'' $\to$ \texttt{[(i,love), (love,cats)]}
    \item \textbf{Trigrammi} ($N=3$): triplette consecutive.
          Es.: ``i love cats'' $\to$ \texttt{[(i,love,cats)]}
\end{itemize}

\subsection{Bigrammi senza stopwords}

La rimozione delle stopwords fa emergere le coppie di parole
più caratteristiche di ciascun dominio.

\begin{table}[H]
\centering
\begin{tabular}{llr}
\toprule
\textbf{Dominio} & \textbf{Bigramma} & \textbf{Frequenza} \\
\midrule
Twitter & (wan, na)          & 123 \\
Twitter & (follow, please)   & 81  \\
Twitter & (gon, na)          & 72  \\
Twitter & (follow, u)        & 62  \\
Twitter & (happy, birthday)  & 54  \\
\midrule
Letterario & (mr, knightley)    & 248 \\
Letterario & (mrs, weston)      & 222 \\
Letterario & (mr, elton)        & 184 \\
Letterario & (miss, woodhouse)  & 148 \\
Letterario & (frank, churchill) & 123 \\
\bottomrule
\end{tabular}
\caption{Top bigrammi per dominio dopo la rimozione delle stopwords}
\end{table}

I bigrammi Twitter mostrano chiaramente le interazioni sociali tipiche
della piattaforma: \textit{wanna} e \textit{gonna} (spezzati dal
tokenizzatore in \texttt{wan na} e \texttt{gon na}), richieste di
follow reciproco, auguri di compleanno. I bigrammi letterari sono
quasi esclusivamente titoli e cognomi dei personaggi del romanzo,
riflettendo la struttura narrativa del testo e il registro formale
dell'inglese ottocentesco.

\subsection{Bigrammi con stopwords}

\begin{table}[H]
\centering
\begin{tabular}{llr}
\toprule
\textbf{Dominio} & \textbf{Bigramma} & \textbf{Frequenza} \\
\midrule
Twitter & (thank, you) & 260 \\
Twitter & (i, love)    & 173 \\
Twitter & (i, miss)    & 160 \\
Twitter & (love, you)  & 155 \\
Twitter & (i, cant)    & 155 \\
\midrule
Letterario & (to, be)    & 602 \\
Letterario & (of, the)   & 563 \\
Letterario & (i, am)     & 363 \\
Letterario & (she, had)  & 322 \\
Letterario & (had, been) & 305 \\
\bottomrule
\end{tabular}
\caption{Top bigrammi per dominio con stopwords}
\end{table}

Alcune osservazioni interessanti:
\begin{itemize}[leftmargin=1.5cm]
    \item Twitter usa \textit{cant} (contrazione senza apostrofo),
          mentre il romanzo usa \textit{i am} (forma estesa): è una
          differenza di registro molto netta.
    \item La struttura \textit{had been} (past perfect) è molto
          frequente nel romanzo (305 occorrenze) e praticamente
          assente nei tweet, dove i tempi verbali sono più semplici.
    \item Twitter è dominato da espressioni affettive dirette
          (\textit{i love}, \textit{love you}, \textit{i miss}),
          tipiche della comunicazione personale sui social.
\end{itemize}

\subsection{Trigrammi}

\begin{table}[H]
\centering
\begin{tabular}{llr}
\toprule
\textbf{Dominio} & \textbf{Trigramma} & \textbf{Frequenza} \\
\midrule
Twitter & (i, love, you)     & 84 \\
Twitter & (you, so, much)    & 74 \\
Twitter & (thanks, for, the) & 67 \\
Twitter & (i, want, to)      & 63 \\
\midrule
Letterario & (i, do, not)      & 123 \\
Letterario & (i, am, sure)     & 97  \\
Letterario & (she, could, not) & 68  \\
Letterario & (a, great, deal)  & 63  \\
\bottomrule
\end{tabular}
\caption{Top trigrammi per dominio}
\end{table}

I trigrammi letterari mostrano costruzioni tipiche del parlato
ottocentesco trascritto: \textit{i do not}, \textit{i am sure},
\textit{she could not}. Twitter al contrario produce espressioni
emotive brevi e ripetitive: \textit{i love you}, \textit{you so much}.

% ============================================================
\section{Modelli di Linguaggio}
% ============================================================

\subsection{Maximum Likelihood Estimation (MLE)}

Il modello MLE stima la probabilità di una parola dato il contesto
precedente semplicemente contando le occorrenze nel corpus di training.
Per i bigrammi, la formula è:

\begin{equation}
P(w_i \mid w_{i-1}) = \frac{C(w_{i-1},\, w_i)}{C(w_{i-1})}
\end{equation}

dove $C(\cdot)$ indica il numero di occorrenze nel corpus.

\begin{lstlisting}[language=Python, caption=Addestramento del modello MLE]
from nltk.lm import MLE
from nltk.lm.preprocessing import padded_everygram_pipeline

n = 2  # bigrammi
train_data, vocab = padded_everygram_pipeline(n, [tokens_tweet])

modello_tweet = MLE(n)
modello_tweet.fit(train_data, vocab)
\end{lstlisting}

La funzione \texttt{padded\_everygram\_pipeline} aggiunge i simboli
\texttt{<s>} (inizio frase) e \texttt{</s>} (fine frase) ai token
e costruisce tutti gli N-gram fino all'ordine $n$ richiesto.

\subsubsection{Probabilità condizionali osservate}

\begin{table}[H]
\centering
\begin{tabular}{llrr}
\toprule
\textbf{Parola} & \textbf{Contesto} &
\textbf{P (Twitter)} & \textbf{P (Letterario)} \\
\midrule
love  & i    & 0,0530 & 0,0017 \\
wait  & cant & 0,1146 & 0,0000 \\
happy & so   & 0,0138 & ---    \\
mr    & said & 0,0000 & 0,0802 \\
been  & had  & ---    & 0,1885 \\
\bottomrule
\end{tabular}
\caption{Probabilità condizionali MLE nei due modelli}
\end{table}

I risultati confermano le aspettative. \textit{love} segue \textit{i}
con probabilità del 5,3\% su Twitter e quasi zero nel romanzo.
\textit{had been} è frequentissimo nel romanzo (18,9\%), riflettendo
l'uso narrativo del past perfect, mentre è praticamente assente nei tweet.

\subsubsection{Il problema dello zero}

Il principale limite del MLE è che assegna probabilità esattamente
\textbf{zero} a qualsiasi sequenza non osservata durante il training.
Questo causa perplexity infinita su qualsiasi testo che contenga
anche una sola parola o coppia mai vista:

\begin{equation}
P_{\text{MLE}}\bigl(\text{``xyzword''} \mid \text{``i''}\bigr)
= \frac{0}{C(\text{``i''})} = 0
\quad \Rightarrow \quad
\text{PP} = \infty
\end{equation}

\begin{lstlisting}[language=Python, caption=Dimostrazione del problema dello zero]
print(modello_tweet.score('xyzword', ['i']))
# -> 0.000000

frase_rara = word_tokenize("i xyzword lol")
print(modello_tweet.perplexity(frase_rara))
# -> inf
\end{lstlisting}

Questo è un problema reale: anche un solo token fuori vocabolario
rende inutilizzabile la perplexity MLE come metrica di valutazione.

\subsection{Smoothing di Laplace}

Il \textbf{Laplace smoothing} (o add-one smoothing) risolve il problema
aggiungendo 1 a tutti i conteggi, incluse le coppie mai osservate:

\begin{equation}
P_{\text{Laplace}}(w_i \mid w_{i-1}) =
\frac{C(w_{i-1},\, w_i) + 1}{C(w_{i-1}) + |V|}
\end{equation}

dove $|V|$ è la dimensione del vocabolario. In questo modo nessuna
sequenza ha probabilità zero, ma le sequenze mai viste ricevono
una probabilità molto piccola.

\begin{lstlisting}[language=Python, caption=Addestramento del modello Laplace]
from nltk.lm import Laplace

modello_lp = Laplace(n)
modello_lp.fit(train_data, vocab)

print(modello_lp.score('xyzword', ['i']))
# -> 0.000065  (piccolo ma non zero)

print(modello_lp.perplexity(frase_rara))
# -> 1677.95  (alta ma non infinita)
\end{lstlisting}

\begin{table}[H]
\centering
\begin{tabular}{lrr}
\toprule
\textbf{Metrica} & \textbf{MLE} & \textbf{Laplace} \\
\midrule
$P(\text{``xyzword''} \mid \text{``i''})$ & 0,000000 & 0,000065 \\
Perplexity su parola mai vista            & $\infty$ & 1677,95  \\
Gestisce OOV\footnotemark                 & No       & Sì       \\
\bottomrule
\end{tabular}
\caption{Confronto MLE vs Laplace}
\end{table}
\footnotetext{OOV = Out-Of-Vocabulary, parole non presenti nel training set.}

\subsection{Scelta del valore di N}

\begin{table}[H]
\centering
\begin{tabular}{lrr}
\toprule
$N$ & \textbf{PP frase Twitter} & \textbf{PP frase Letterario} \\
\midrule
1 & 5170,14 & 12025,00 \\
2 & 5170,91 & 12027,00 \\
3 & 5170,92 & 12027,00 \\
\bottomrule
\end{tabular}
\caption{Perplexity Laplace per $N=1,2,3$ (modello Twitter)}
\end{table}

I valori quasi identici per $N=1,2,3$ sono sintomo di un problema
di \textbf{data sparsity}: con soli 94.706 token Twitter, la
maggior parte dei bigrammi e trigrammi possibili appare raramente
o mai nel corpus. Laplace tratta quindi quasi tutte queste sequenze
come sconosciute, producendo perplexity simili indipendentemente da $N$.
Con un corpus più grande le differenze sarebbero più marcate, e i
trigrammi catturerebbero strutture sintattiche più complesse.

% ============================================================
\section{Perplexity e Classificazione}
% ============================================================

\subsection{Perplexity}

La \textbf{perplexity} di un modello su una sequenza $W = w_1 \ldots w_N$
misura quanto il modello è ``sorpreso'' da quel testo:

\begin{equation}
\text{PP}(W) = P(w_1, \ldots, w_N)^{-\frac{1}{N}}
= \sqrt[N]{\prod_{i=1}^{N} \frac{1}{P(w_i \mid w_{i-1})}}
\end{equation}

Un valore \textbf{basso} significa che il modello assegna alta probabilità
alla sequenza (la conosce bene); un valore \textbf{alto} indica sorpresa.

\subsection{Limiti della perplexity per la classificazione}

\begin{table}[H]
\centering
\begin{tabular}{llrrl}
\toprule
\textbf{Frase} & \textbf{Atteso} &
\textbf{PP Twitter} & \textbf{PP Lett.} & \textbf{Risultato} \\
\midrule
``i love you so much''     & TW & 5170,9  & 4488,5  & TW \checkmark \\
``so happy right now''     & TW & 12027,0 & 9308,0  & LT $\times$   \\
``i miss you feel sad''    & TW & 5746,3  & 4917,0  & TW \checkmark \\
``she had been very well'' & LT & 12027,0 & 9308,0  & TW $\times$   \\
``it was a very good''     & LT & 6756,2  & 5258,0  & TW $\times$   \\
\bottomrule
\end{tabular}
\caption{Classificazione con perplexity Laplace: accuratezza 3/5}
\end{table}

Il classificatore basato su perplexity è inaffidabile. Il motivo è
strutturale: il modello Twitter ha un vocabolario più grande (12.027
parole) rispetto a quello letterario (9.308 parole). Con lo smoothing
di Laplace, un vocabolario più grande ``diluisce'' la probabilità su
più termini, producendo sistematicamente perplexity più alta.
Il classificatore quindi tende sempre a preferire il modello con
vocabolario più piccolo, indipendentemente dal contenuto del testo.

\subsection{Classificatore con Log-Likelihood}

Per superare questo limite si è adottata la \textbf{log-likelihood}
come metrica di classificazione:

\begin{equation}
\mathcal{L}(W) = \log P(w_1, \ldots, w_N)
= \sum_{i=2}^{N} \log P(w_i \mid w_{i-1})
\end{equation}

Un valore \textbf{meno negativo} (più vicino a zero) indica che il
modello assegna alta probabilità alla sequenza, ovvero la ``conosce bene''.
Poiché la log-likelihood non dipende dalla dimensione del vocabolario
in modo diretto come la perplexity, è una metrica più robusta per
confrontare modelli su vocabolari di dimensioni diverse.

\begin{lstlisting}[language=Python, caption=Implementazione del classificatore log-likelihood]
import math

def log_likelihood(modello, tokens):
    score = 0.0
    for i in range(1, len(tokens)):
        prob = modello.score(tokens[i], [tokens[i-1]])
        score += math.log(prob) if prob > 0 else math.log(1e-10)
    return score

def classifica(testo, mod_tw, mod_lt):
    tokens = word_tokenize(testo.lower())
    ll_tw   = log_likelihood(mod_tw, tokens)
    ll_lett = log_likelihood(mod_lt, tokens)
    return "TWITTER" if ll_tw > ll_lett else "LETTERARIO"
\end{lstlisting}

\subsubsection{Risultati}

\begin{table}[H]
\centering
\begin{tabular}{p{7.5cm}rrr}
\toprule
\textbf{Testo} & \textbf{LL Twitter} & \textbf{LL Lett.} & \textbf{Verdetto} \\
\midrule
``i love you so much thank you''
    & $-32,02$ & $-42,15$ & TWITTER \checkmark \\
``she had been very well disposed''
    & $-124,74$ & $-21,91$ & LETTERARIO \checkmark \\
``i miss you so much cant stop thinking''
    & $-41,54$ & $-122,88$ & TWITTER \checkmark \\
``it was a very good thing to have done''
    & $-29,21$ & $-24,75$ & LETTERARIO \checkmark \\
``happy birthday love you so much''
    & $-11,27$ & $-58,54$ & TWITTER \checkmark \\
``she was not a woman of many words''
    & $-70,82$ & $-47,21$ & LETTERARIO \checkmark \\
``i cant stop smiling so happy''
    & $-57,95$ & $-97,56$ & TWITTER \checkmark \\
``mr knightley had always been very fond''
    & $-108,42$ & $-24,91$ & LETTERARIO \checkmark \\
\bottomrule
\end{tabular}
\caption{Risultati classificatore log-likelihood: 8/8 corretti}
\end{table}

Il classificatore log-likelihood classifica correttamente tutti gli
8 casi di test, rispetto ai 3/8 del classificatore basato su perplexity.
Le differenze di log-likelihood sono particolarmente marcate per testi
fortemente tipici di un dominio: ``she had been very well disposed''
ha LL letterario di $-21,91$ contro $-124,74$ per Twitter, una
differenza di oltre 100 unità che riflette quanto quella struttura
sintattica sia rara nei tweet.

% ============================================================
\section{Generazione del Testo}
% ============================================================

I modelli MLE possono essere usati per \textbf{generare} testo nello
stile del corpus su cui sono stati addestrati. La generazione è
un processo probabilistico: ad ogni passo, la parola successiva
viene campionata dalla distribuzione condizionale data la parola
precedente.

Si è scelto di usare MLE (e non Laplace) per la generazione, perché
Laplace distribuisce la probabilità su tutti i token del vocabolario
inclusi quelli sconosciuti, producendo sequenze di lettere singole
invece di parole reali.

\begin{lstlisting}[language=Python, caption=Generazione del testo]
def genera_testo(modello, n_parole=40):
    testo_raw = modello.generate(n_parole, random_seed=42)
    # Filtra token speciali e artefatti
    testo_filtrato = [
        t for t in testo_raw
        if t not in ['<s>', '</s>', '<UNK>']
        and not (len(t) == 1 and t not in ['i', 'a'])
    ]
    return ' '.join(testo_filtrato)
\end{lstlisting}

\begin{table}[H]
\centering
\begin{tabular}{p{2.2cm}p{10cm}}
\toprule
\textbf{Stile} & \textbf{Testo generato (seed=42)} \\
\midrule
Twitter &
``papa coach it from the rain why are not a fantastic spring
fiesta make sure the donation manager we all the rest his cute
when im being a small polka dots'' \\[0.5cm]
Letterario &
``over a good exchange consent to the cardroom looking about
her less allied with persevering attention could not that a
substance of her but yet heard any body to me'' \\
\bottomrule
\end{tabular}
\caption{Esempi di testo generato con i due modelli (seed=42)}
\end{table}

Il testo Twitter è frammentato e colloquiale, con espressioni
informali (\textit{im}, \textit{polka dots}). Il testo letterario
usa strutture più elaborate (\textit{persevering attention},
\textit{allied with}) tipiche della prosa ottocentesca.

\begin{table}[H]
\centering
\begin{tabular}{cp{5.5cm}p{5.5cm}}
\toprule
\textbf{Seed} & \textbf{Twitter} & \textbf{Letterario} \\
\midrule
99  & ``im getting bigger at the epic i love you in my life going'' &
      ``i can believe i shall decidedly in the neighbourhood who'' \\
7   & ``have a new cards in march free insurance co i beleve'' &
      ``hartfield and return among their getting at last ah my'' \\
123 & ``and by my bestfriend she brought us i want corn chips'' &
      ``and away from being to announce cheerful exulting than'' \\
\bottomrule
\end{tabular}
\caption{Generazione con diversi seed casuali}
\end{table}

% ============================================================
\section{Analisi Linguistica}
% ============================================================

\subsection{Type-Token Ratio}

Il \textbf{Type-Token Ratio} (TTR) misura la ricchezza lessicale
di un corpus come rapporto tra tipi (parole distinte) e token
(parole totali):

\begin{equation}
\text{TTR} =
\frac{|\{w : w \in \text{corpus}\}|}{|\text{corpus}|}
= \frac{\text{types}}{\text{tokens}}
\end{equation}

\begin{table}[H]
\centering
\begin{tabular}{lrrl}
\toprule
\textbf{Corpus} & \textbf{Types} & \textbf{Tokens} & \textbf{TTR} \\
\midrule
Twitter    & 12.024 & 94.706  & 0,1270 \\
Letterario & 9.305  & 158.267 & 0,0588 \\
\bottomrule
\end{tabular}
\caption{Type-Token Ratio dei due corpora}
\end{table}

Il TTR di Twitter è più del doppio rispetto al romanzo.
Questo risultato è controintuitivo: ci si aspetterebbe che un romanzo
ottocentesco abbia un vocabolario più ricco di tweet. La spiegazione
sta nella natura eterogenea dei tweet: ogni tweet introduce hashtag,
abbreviazioni, slang e nomi propri unici (\texttt{followfriday},
\texttt{unfollowers}, \texttt{bestfriend}), molti dei quali appaiono
raramente nel corpus. Il romanzo invece ripete sistematicamente gli
stessi personaggi e costruzioni grammaticali, abbassando il rapporto
types/tokens.

\subsection{Lunghezza media delle frasi}

\begin{table}[H]
\centering
\begin{tabular}{lr}
\toprule
\textbf{Corpus} & \textbf{Lunghezza media (parole)} \\
\midrule
Tweet             & 16,0 \\
Frase letteraria  & 25,6 \\
\bottomrule
\end{tabular}
\caption{Lunghezza media delle frasi nei due corpora}
\end{table}

Le frasi letterarie sono in media 9,6 parole più lunghe dei tweet.
Questo riflette la struttura sintattica complessa del romanzo
ottocentesco, con proposizioni subordinate multiple, incisi e
descrizioni elaborate, contrapposta alla brevità imposta dal
limite di 140 caratteri di Twitter (al momento della raccolta
del dataset nel 2015).

% ============================================================
\section{Limiti, Vantaggi e Confronto}
% ============================================================

\subsection{Vantaggi degli N-gram}

\begin{itemize}[leftmargin=1.5cm]
    \item \textbf{Semplicità e trasparenza}: l'intero modello si riduce
          a conteggi di sequenze. Le probabilità condizionali sono
          direttamente leggibili e interpretabili.
    \item \textbf{Efficienza computazionale}: addestramento e inferenza
          sono molto veloci, non richiedono GPU né risorse elevate.
    \item \textbf{Nessun training pre-esistente richiesto}: si addestrare
          da zero su qualsiasi corpus in pochi secondi.
\end{itemize}

\subsection{Limiti degli N-gram}

\begin{itemize}[leftmargin=1.5cm]
    \item \textbf{Data sparsity}: con dataset piccoli, la maggior parte
          dei bigrammi e trigrammi possibili non viene mai osservata nel
          training. Nel nostro caso, i 94.706 token Twitter sono
          insufficienti per stimare bene le probabilità di trigrammi,
          come dimostrano i valori di perplexity quasi identici per
          $N=1,2,3$.

    \item \textbf{Contesto limitato}: i bigrammi considerano solo la
          parola immediatamente precedente. Non possono catturare
          dipendenze a lungo raggio come in ``The man who I met
          yesterday \textbf{was} kind'', dove \textit{was} dipende
          da \textit{man} che è lontano.

    \item \textbf{Nessuna semantica}: il modello non capisce il
          significato delle parole. \textit{happy} e \textit{joyful}
          sono token completamente diversi, anche se sinonimi.

    \item \textbf{Perplexity non comparabile tra vocabolari diversi}:
          come dimostrato nella sezione sulla classificazione, la
          perplexity con Laplace non è direttamente comparabile tra
          modelli con vocabolari di dimensioni diverse. È stato
          necessario ricorrere alla log-likelihood per ottenere un
          classificatore affidabile.
\end{itemize}

\subsection{Confronto con approcci alternativi}

\begin{table}[H]
\centering
\begin{tabular}{p{2.5cm}p{3.5cm}p{3.5cm}p{2cm}}
\toprule
\textbf{Approccio} & \textbf{Vantaggi} &
\textbf{Svantaggi} & \textbf{Complessità} \\
\midrule
N-gram &
Semplice, veloce, interpretabile &
Sparsità, no semantica, contesto corto &
Bassa \\
\midrule
Word2Vec &
Semantica, vettori densi &
No probabilità, no generazione &
Media \\
\midrule
LSTM / RNN &
Contesto lungo, genera testo &
Richiede molto training &
Alta \\
\midrule
Transformer &
Stato dell'arte &
Computazionalmente costoso &
Molto alta \\
\bottomrule
\end{tabular}
\caption{Confronto tra approcci per la modellazione del linguaggio}
\end{table}

% ============================================================
\section{Risultati e Conclusioni}
% ============================================================

\subsection{Riepilogo}

\begin{table}[H]
\centering
\begin{tabular}{p{4cm}p{4cm}p{5cm}}
\toprule
\textbf{Esperimento} & \textbf{Risultato} & \textbf{Interpretazione} \\
\midrule
Bigrammi Twitter     & (wan,na), (i,love), (i,cant) &
Slang, espressioni emotive \\
Bigrammi Letterario  & (mr,knightley), (had,been) &
Titoli formali, past perfect \\
TTR Twitter          & 0,1270 &
Vocabolario eterogeneo \\
TTR Letterario       & 0,0588 &
Vocabolario ripetitivo \\
Lunghezza frasi TW   & 16,0 parole &
Brevi, imposte dal limite 140 car. \\
Lunghezza frasi LT   & 25,6 parole &
Complesse, subordinate multiple \\
Classif. perplexity  & 3/8 = 38\% &
Inaffidabile con vocabolari diversi \\
Classif. log-lik.    & 8/8 = 100\% &
Robusto, non dipende da $|V|$ \\
\bottomrule
\end{tabular}
\caption{Riepilogo dei principali risultati}
\end{table}

\subsection{Conclusioni}

Questo laboratorio ha mostrato che i modelli N-gram, pur nella loro
semplicità, sono strumenti efficaci per caratterizzare e distinguere
stili linguistici diversi.

I principali risultati ottenuti sono:

\begin{enumerate}[leftmargin=1.5cm]
    \item I bigrammi e trigrammi più frequenti rispecchiano fedelmente
          le caratteristiche culturali e stilistiche di ciascun dominio:
          slang e interazioni sociali per Twitter, titoli formali e
          strutture narrative complesse per il romanzo.

    \item Il TTR più alto di Twitter è controintuitivo ma spiegabile:
          la varietà lessicale dei tweet deriva dall'eterogeneità dei
          contenuti e dall'uso di hashtag e abbreviazioni uniche,
          non da un vocabolario intrinsecamente più ricco.

    \item La perplexity con Laplace non è una metrica affidabile per
          classificare testi tra modelli con vocabolari di dimensioni
          diverse. La log-likelihood è una metrica più robusta
          per questo scopo, raggiungendo il 100\% di accuratezza
          sui casi di test.

    \item I limiti osservati (sparsità, perplexity identica per $N=1,2,3$)
          sono in gran parte attribuibili alla dimensione ridotta del
          corpus Twitter. Un dataset più grande ridurrebbe la sparsità
          e farebbe emergere differenze più nette tra unigrammi,
          bigrammi e trigrammi.
\end{enumerate}

Sviluppi futuri potrebbero includere: l'uso di dataset più grandi,
l'adozione di smoothing più sofisticati come Kneser-Ney, e il
confronto con modelli neurali (LSTM, Transformer) per valutare
il guadagno in qualità rispetto alla complessità computazionale.

% ============================================================
\section*{Dichiarazione sull'uso di strumenti generativi}
% ============================================================

\addcontentsline{toc}{section}{Dichiarazione sull'uso di strumenti generativi}

Come richiesto dalle linee guida del laboratorio, si dichiara che
nella realizzazione di questo lavoro sono stati utilizzati strumenti
di intelligenza artificiale generativa (Claude, Anthropic) per:

\begin{itemize}[leftmargin=1.5cm]
    \item Supporto nella strutturazione e nel debugging del codice Python.
    \item Identificazione e correzione di errori nel classificatore
          e nella generazione del testo.
    \item Formattazione della presente relazione in \LaTeX.
\end{itemize}

Tutte le scelte progettuali, l'analisi e l'interpretazione dei
risultati, e le conclusioni sono state elaborate autonomamente
dal gruppo.

% ============================================================
\end{document}
% ============================================================